\documentclass[../main.tex]{subfiles}

\begin{document}

\chapter*{\myAbstractTitle}

In a federated identity infrastructure, a single identity provider authenticates users on behalf 
of many service providers and thereby becomes the trust anchor of the entire federation. 
A compromised or misconfigured provider does not affect one application but undermines the 
security of every service relying on it. 
Despite this exposure, systematic security monitoring at the application layer of such providers 
remains largely absent, and existing approaches in this domain focus on operational availability 
rather than on security.

This thesis presents a monitoring system that observes a Shibboleth identity provider from within 
its own application layer. 
The system draws on three sources that differ in what they can observe, namely the audit log 
written at the end of a transaction, the configuration and metadata declared for each service 
provider, and the runtime state accessible through an interceptor executing inside the 
authentication flow. 
Records from these sources are pseudonymised at their origin, transmitted to a backend located 
outside the trust boundary of the provider, correlated with one another and evaluated against a 
catalogue of thirteen detection rules. 
The rules are not assembled from experience but derived from vulnerabilities and attack classes 
documented in a systematic survey of federated identity security and from the normative 
requirements of the SAML security considerations.

The evaluation addresses three questions, namely which data can be collected, how reliably it can 
be correlated, and what the resulting system costs. 
The correlation reaches a coverage of 99,9\% and remains free of incorrect assignments up to a 
quantified login density, beyond which the heuristic strategies can no longer resolve the order of 
events. 
All thirteen rules were verified against a real search index and five additionally through the 
complete pipeline from a browser login to the resulting alert. 
In the productive transmission mode, no overhead was measurable in the end-to-end login latency, 
with the work performed inside the interceptor amounting to approximately 0,3\% of the duration of 
a login.

The work also identifies the boundaries of the chosen observation point. 
Attack classes located at a different system layer, involving a different party, or observable 
only through their preconditions remain outside its reach, and a provider abusing its own position 
cannot be detected by a system executing inside it. 
Beyond these results, the evaluation uncovered three defects in the monitoring infrastructure that 
operation itself had not revealed, which illustrates the premise underlying the work, namely that 
a system appearing to function is not thereby known to function.

\end{document}
